\section{引言}
\label{sec:intro}

行人轨迹预测以场景中行人近期的位置序列为输入，推断其在未来数秒内的运动位置，通常以多条候选轨迹及其概率的形式输出。它是自动驾驶、移动机器人导航与人群安全管理的关键环节，预测结果直接进入下游规划器，用于判断其他智能体的运动意图。过去十年，该方向在两个方面持续推进：一是社会交互建模，从对邻居隐状态做社会池化的 Social LSTM 与 Social GAN \cite{alahi2016social,gupta2018social}，发展到基于图结构与注意力机制的场景编码器 \cite{mohamed2020social,yuan2021agentformer,liang2020lane,gao2020vectornet}；二是未来分布建模，从单一轨迹回归发展到以生成模型刻画未来多模态性的解码器 \cite{mangalam2020not,salzmann2020trajectron,fu2025moflow,zhang2024demo}。

上述研究几乎共享一个很少被明确表述的前提：历史轨迹完整、连续且等间隔采样。输入被视作一个 $T\times 2$ 的张量，其中每一行都是一次测量。真实系统并不提供这样的张量。历史轨迹由检测器与多目标跟踪器拼接而成，并继承该流程的全部失效模式：行人走到停放的车辆后方，检测在连续四帧中丢失；两人交错而行，跟踪器交换了身份；计算负载升高时相机丢帧。其结果是，每个行人、每个样本中\emph{被观测到的帧集合}各不相同，而缺失行不携带任何位置信息。

面对不完整历史，最常见的做法是先把历史补全再预测：由统计或学习型插补器填补空缺，再把修复后的序列送入现成的预测器 \cite{chib2024pedestrian,xu2023uncovering}；较新的工作则将两个阶段耦合、联合优化 \cite{chib2024mstip,chib2026hypergraph,sun2025rmpyolo}。两类设计都存在结构性问题。插补以观测数据分布上的重建误差为目标，而预测以未来位移误差为评价标准；只有当重建值等于真实位置的条件均值时，两个目标才一致，而这在缺失较长或运动多模态时恰恰不成立。更严重的是，缺口被填补之后，预测器无法再区分哪些行是测量、哪些行是推断，一个被“编造”出来的坐标会以与测量同等的地位进入速度估计、进入与邻居的社会关系计算，最终进入预测结果。

另一条路线是让预测器直接消费不完整历史，并把每一帧可用的\emph{证据}建模为显式量。这一视角带来三点与插补视角不同的推论。第一，历史状态的可靠程度不是单帧属性，而是序列属性：它随最近一次观测的年龄衰减，并应调节该状态进入时间递推的强度。第二，缺失性改变了社会交互的几何结构：当前未被观测的行人仍可以查询邻居，但其陈旧位置不应作为当前证据广播给其他行人。第三，解码起点本身不确定：当末端若干帧缺失时，最近的可信位置早于预测起点，解码器需要被告知这一间隔，而不是被锚定在一个占位行上。

实现这一思路需要两样目前缺失的东西。其一是\emph{协议}。已发表的不完整历史研究在缺失生成方式、样本划分、检查点选择与排序指标上并不一致；部分工作在同一个表格中并列报告“先插补后预测”与“直接预测”，而两种设置向模型暴露的信息并不相同。因此跨论文的数值不可比，单一设计选择的作用也无法被隔离。其二是贯穿全模型的\emph{机制}。现有缺失感知方法多把二值掩码拼接为输入通道 \cite{chen2023trajmae,yang2023rmp,chib2023dropping}，或将掩码用作预训练目标，或只考虑历史被整体截短的情形 \cite{sun2022momentary,monti2022howmany,qiu2025adapting}；它们都没有把证据状态作为一等的变量，用以同时决定哪些行人可以作为空间证据、每一帧以多大强度进入时间递推、以及解码器从哪里开始。

本文针对第三视角（固定摄像机或鸟瞰视角）场景回应上述两个缺口。首先冻结 \textbf{TrajGap-Bench} 评测基准。该基准建立在公开的 TrajImpute 数据发布之上 \cite{chib2024pedestrian}，其贡献不是重新发布轨迹，而是固定实验契约：将 Easy 与 Hard 两个版本在样本级合并为统一的 Mixed 划分，用于训练、验证与测试；对所有方法施加同一条基于最终预测头验证集 $\mathrm{minFDE}_{20}$ 的检查点规则；仅以 $K=20$ 的 $\mathrm{minADE}_{20}$ 与 $\mathrm{minFDE}_{20}$ 两项指标定义排序；并将缺失帧数、锚点间隔、最长连续缺失、有效帧数、末端缺失五个维度作为分层诊断而非额外榜单报告。随后提出 \textbf{GAST}（Gap-Aware Spatio-Temporal forecaster）。该框架以观测证据状态表示每个行人的历史，用其门控逐帧关系编码器、调制因果状态空间时间编码器及其读出，并以记录“末次观测距预测起点多远”的锚点状态条件化目标端点式多模态解码器。模型不对任何坐标做插补；当历史完整时，证据条件化退化为恒等映射，模型回落到其完整历史骨干。

本文在基准的全部五个场景、共 302{,}886 个测试样本上，将 GAST 与六种外部基线对比。主要发现如下。
\begin{itemize}
\item GAST 取得所有对比方法中最优的宏平均 $\mathrm{minADE}_{20}$/$\mathrm{minFDE}_{20}$（0.3168/0.4831~m），相对最强外部基线（MoFlow）降低 9.4\%/7.6\%；对 MoFlow 在 $\mathrm{minADE}_{20}$ 上五场景全胜，在 $\mathrm{minFDE}_{20}$ 上四胜一负。
\item 优势随缺失程度扩大而非均匀分布。相对 MoFlow，$\mathrm{minFDE}_{20}$ 在历史完整时降低 5.4\%，缺失三帧时降低 11.2\%，六帧时降低 22.6\%；而在仅剩一个观测的极端设定下反转为落后约 4\%。这一“剂量—效应”形态与证据条件化的设计意图一致。
\item 本文同时报告两项限定性结果：其一，改进并不跨场景一致（GAST 在 HOTEL-M 上落后于 AgentFormer，在 ZARA1-M 的 $\mathrm{minFDE}_{20}$ 上落后于 MoFlow）；其二，best-of-20 的提升未同步体现在概率最高的单一模式上，top-1 指标与 MoFlow 基本相当。本文对二者进行分析，并指出它们如何约束该方法的使用与扩展。
\end{itemize}

本文余下部分安排如下：第~\ref{sec:related} 节综述相关工作；第~\ref{sec:bench} 节给出任务定义与评测协议；第~\ref{sec:method} 节介绍 GAST；第~\ref{sec:exp} 节报告实验与分层分析；第~\ref{sec:disc} 节讨论结果及其局限；第~\ref{sec:concl} 节总结。
